\chapter{Avant-propos}

\section{Pourquoi cet atelier existe}

Les règles de Conqueror ne sont pas connues. Plusieurs sont ouvertement
suspectes, et aucune ne se tranche par le débat. La version précédente du
document de règles s'annonçait « TOTALEMENT VERROUILLÉE » avec « 16 décisions
prises » ; plusieurs de ces décisions étaient mesurablement mauvaises, et
personne ne pouvait le savoir sans les simuler.

D'où la règle de travail du projet, qui commande tout le reste.

\begin{nkcode}{REGLES\_COMPLETES\_v2.md:31-35}
La regle de travail de ce projet. Toute valeur numerique de ce document est
un PARAMETRE NOMME, jamais une constante. Le moteur de regles doit exposer
ces parametres pour qu'ils soient modifiables SANS RECOMPILATION. Une regle
n'est pas defendue par argument : elle est refutee ou confirmee par
simulation de masse.
\end{nkcode}

ConquerorLab est l'instrument qui produit cette simulation de masse. Il permet
de \textbf{tester une règle en dix minutes au lieu de trois semaines}.

\section{À qui s'adresse ce cours}

\begin{center}
\begin{tabular}{@{}p{6.4cm}p{7.4cm}@{}}
\toprule
\textbf{Vous êtes\ldots} & \textbf{Lisez au minimum} \\
\midrule
\textbf{A1} — vous écrivez le moteur de règles & chapitres 0, 1, 2, 4, 5, 6 \\
\textbf{A2} — vous écrivez l'IA adversaire & chapitres 0, 1, 3, 6 \\
\textbf{designer} — vous réglez et vous mesurez & chapitres 0, 1, 4, 6 \\
\textbf{vous reprenez le projet} & tout, dans l'ordre \\
\bottomrule
\end{tabular}
\end{center}

Prérequis : savoir écrire du C++ — pas du C++ avancé. Les trois exemples de ce
cours n'utilisent ni templates, ni héritage, ni exceptions. Aucune connaissance
de Nkentseu n'est nécessaire pour les chapitres 2, 3 et 5.

\section{Le kit, et sur quoi il tourne}

Vous avez reçu un dossier — ou une archive \texttt{ConquerorLab-Kit.zip} — qui
contient tout : l'atelier, les en-têtes, les bibliothèques, les exemples, et ce
cours en PDF. \textbf{Il n'y a rien à installer sauf un compilateur} (voir
\texttt{LISEZMOI.txt}).

\begin{nkcode}{Disposition du kit}
ConquerorLab.exe          l'atelier
include/                  TOUS les en-tetes, a plat
    Conqueror/            le contrat
    NKCore/  NKMath/  NKContainers/  NKLogger/  ...
lib/                      les bibliotheques liees a vos modules
travail/rules/            <- vos moteurs de regles    (.cpp)
travail/ai/               <- vos IA                   (.cpp)
travail/boards/           <- vos grilles              (.json)
exemples/                 trois modules complets, a recopier
\end{nkcode}

\begin{nkretenir}[\texttt{include/} reflète exactement ce que vous tapez]
\nkfile{include/NKMath/NkFunctions.h} correspond à
\texttt{\#include "NKMath/NkFunctions.h"}.

La première version du kit recopiait l'arborescence du dépôt —
\nkfile{Kernel/Foundation/NKCore/src/NKCore/...}. \texttt{Foundation},
\texttt{System}, \texttt{src} sont de la plomberie de dépôt : trois niveaux qui
ne veulent rien dire quand on écrit des règles de jeu, et qui séparent les
modules de ce qu'ils incluent. L'atelier n'a donc qu'\textbf{un seul \texttt{-I}}.
\end{nkretenir}

\begin{nkpiege}[Windows uniquement, pour l'instant]
Le kit livré est \textbf{Windows x64}. Ce n'est pas une limite de conception :
l'atelier est écrit sur NKGui/NKCanvas, dont les portages Linux, Android et Web
existent. C'est une limite de \textbf{ce qui a été construit et vérifié à ce
jour} — et livrer un binaire non testé serait pire que ne rien livrer.

Si vous travaillez sous Linux ou macOS, vous pouvez cloner le dépôt et compiler
vous-même. Dites-le : c'est le genre de retour qui fait avancer la liste.
\end{nkpiege}

\subsection{Compiler pour une autre plateforme}

Le dépôt se compile avec \texttt{jenga}, l'outil de construction de la maison :

\begin{nkcode}{Compilation depuis les sources}
git clone <depot Nkentseu>
cd Nkentseu
jenga build --target ConquerorLab --config Release
\end{nkcode}

Le binaire sort dans \texttt{Build/Bin/Release-<Plateforme>/ConquerorLab/}.

Deux choses à savoir avant de vous lancer :

\begin{itemize}
  \item \textbf{Le compilateur de modules est PC seulement.} Sur Android et Web
        il n'y a pas de compilateur sur l'appareil : seuls les modules compilés
        \emph{dans} l'application sont disponibles. L'atelier reste jouable et
        mesurable ; c'est le rechargement à chaud qui disparaît. Limite réelle et
        assumée.
  \item \textbf{Les chemins des bibliothèques changent.} \nkfile{LibDir()} vise
        \nkfile{Build/Lib/<Config>-<Plateforme>}, et l'extension des modules
        passe de \texttt{.dll} à \texttt{.so} (\texttt{.dylib} sur macOS). Tout
        cela est déjà écrit et compilé conditionnellement — mais jamais
        \textbf{exécuté} ailleurs que sur Windows. Attendez-vous à corriger deux
        ou trois choses, et signalez-les.
\end{itemize}

\subsection{Le kit évoluera}

Ce n'est pas une livraison figée. Les retours — un message d'erreur
incompréhensible, un panneau qui déborde, une fonction du contrat qui manque —
donnent lieu à une nouvelle version du kit et du cours.

\begin{nkretenir}
Les deux nouveautés les plus récentes sont nées exactement comme ça : le panneau
\emph{Sortie} et l'ouverture de la géométrie au C++ (chapitre 5) viennent tous
deux d'une remarque, pas d'un plan.
\end{nkretenir}

Le fichier \texttt{VERSION.txt} du kit dit quelle version vous avez. Citez-la
dans vos retours.

\section{Lancer l'atelier}

\begin{nkcode}{Applications/ConquerorLab/README.md — section « Lancer »}
cd D:\Projets\2026\Nkentseu\Nkentseu
jenga build --target ConquerorLab --config Release
.\Build\Bin\Release-Windows\ConquerorLab\ConquerorLab.exe
\end{nkcode}

Trois dossiers sont créés au premier lancement :

\begin{nkcode}{Les trois dossiers de contribution}
depot :   Build/ConquerorLab/rules|ai|boards
kit   :   travail/rules|ai|boards
\end{nkcode}

Le chemin diffère parce que « Build » est un mot de dépôt : dans un kit, personne
ne construit rien. C'est \nkfile{NkcLayout.h} qui détecte la disposition, et lui
seul les connaît toutes les deux. \textbf{En cas de doute, ne devinez pas :} le
panneau \emph{Règles} affiche en clair le chemin exact où il regarde.

\textbf{\texttt{boards}, au pluriel, et \texttt{.json} à la fin.} Un dossier
nommé \texttt{board} ou un fichier sans extension ne sont pas lus. Depuis le
2026-08-23 l'atelier vous le \textbf{dit} --- il nomme les fichiers qu'il ignore
et le dossier voisin mal nommé --- mais il ne devine pas à votre place.

Le troisième contient déjà \nkfile{exemple\_plateau\_par\_defaut.json}, écrit par
l'atelier : c'est le plateau du moteur de référence, exporté. Un format qu'on
peut lire vaut mieux qu'un format qu'on décrit.

\section{La première minute}

À l'ouverture, \textbf{rien ne bouge}, et c'est voulu :

\begin{itemize}
  \item le \textbf{joueur 0 est humain}, les autres sont pilotés par l'IA de
        référence ;
  \item la partie est \textbf{en pause} : c'est vous qui la lancez.
\end{itemize}

La barre du plateau dit ce que l'atelier attend. Elle affiche à tout instant
l'une de ces phrases :

\begin{center}
\begin{tabular}{@{}p{6.6cm}p{7.2cm}@{}}
\toprule
\textbf{Ce que vous lisez} & \textbf{Ce que ça veut dire} \\
\midrule
au tour du joueur humain & à vous : cliquez un totem, puis une case verte \\
en pause — Lecture ou Pas à pas & c'est à une IA de jouer, mais rien ne tourne \\
l'IA réfléchit & un thread calcule ; l'interface reste vivante \\
rejeu en pause & vous relisez le passé, le jeu attend \\
partie terminée & le décompte est affiché au centre \\
IA introuvable & un vrai problème : allez au panneau Modules \\
\bottomrule
\end{tabular}
\end{center}

\begin{nkretenir}
\textbf{Une interface immobile et muette se lit comme une panne.} C'est
exactement ce qui s'est passé pendant le développement : le premier réglage
mettait le joueur 0 en humain, l'atelier attendait un clic, et personne ne
comprenait pourquoi « la simulation ne marchait pas ». La barre d'état a été
ajoutée pour cela. \textbf{Si l'atelier semble bloqué, lisez-la avant de
chercher un bug.}
\end{nkretenir}

\section{Trois gestes pour commencer}

\begin{enumerate}
  \item \textbf{Jouer un coup.} Cliquez un de vos totems. Les destinations
        légales s'entourent de vert. Survolez-en une : les totems ennemis
        qu'elle \emph{retournerait} s'entourent de rouge. C'est la lecture
        tactique centrale du jeu — « quelle surface de contact suis-je en train
        d'offrir ? ». Cliquez pour jouer.
  \item \textbf{Laisser tourner.} Bouton \textbf{Lecture}. L'IA répond. Le
        dernier coup est marqué d'un anneau orange qui pulse une seconde.
  \item \textbf{Passer en mode mesure.} Bouton \textbf{IA vs IA} : tous les
        sièges passent à l'IA et la partie tourne seule. C'est la configuration
        dans laquelle l'atelier répond à une question.
\end{enumerate}

\section{Ce que ce cours ne couvre pas}

\begin{itemize}
  \item \textbf{Le jeu lui-même.} Les règles de Conqueror sont dans
        \nkfile{REGLES\_COMPLETES\_v2.md}. Ce cours explique l'outil qui les teste.
  \item \textbf{NKGui, NKCanvas, le moteur.} Un module de stagiaire n'y touche
        pas. Pour modifier l'\emph{interface} de l'atelier, c'est le cours
        \emph{NkCanvas \& NKGui} qu'il vous faut.
  \item \textbf{Les paliers 1 et 2} (fusion, ressources, pouvoirs, artefacts). Le
        moteur de référence livré s'arrête au palier 0. Le contrat, lui, les
        prévoit déjà : c'est pour cela que \texttt{NkcMove} porte des champs
        \texttt{fuseCells} et \texttt{powerId} qui restent à zéro aujourd'hui.
\end{itemize}
